\chapter{Ковариация, корелация}

\section{Съвместна функция на разпределение и~независимост}

\begin{defbox}[~(Съвместна функция на разпределение)]
Нека $X,Y$ са случайни величини. Функцията
\[
  F_{X,Y}(x,y) := \mathbb P(X<x,\,Y<y), \qquad (x,y)\in\mathbb R^2,
\]
се нарича \textbf{съвместна функция на разпределение} на двойката
$(X,Y)$. За нея са в~сила
\[
  F_{X,Y}(\infty,y)=F_Y(y),\qquad F_{X,Y}(x,\infty)=F_X(x),\qquad F_{X,Y}(\infty,\infty)=1,\qquad F_{X,Y}(-\infty,y)=0.
\]
\end{defbox}

\begin{defbox}[~(Независимост чрез съвместната функция на разпределение)]
Случайните величини $X$ и~$Y$ са \textbf{независими}, ако
\[
  F_{X,Y}(x,y) = F_X(x)\cdot F_Y(y), \qquad \forall (x,y)\in\mathbb R^2,
\]
т.е.\ $\mathbb P(X<x,\,Y<y) = \mathbb P(X<x)\cdot\mathbb P(Y<y)$.
\end{defbox}

\begin{corbox}[]
Ако $X,Y$ са дискретни, дефиницията е еквивалентна на
$\mathbb P(X=x_i,Y=y_j) = \mathbb P(X=x_i)\cdot\mathbb P(Y=y_j)$ за
всички $i,j$.
\end{corbox}

\begin{proofbox}
($\Leftarrow$) Нека $\mathbb P(X=x_i,Y=y_j)=\mathbb P(X=x_i)\mathbb
P(Y=y_j)$ за всички $i,j$. Тогава за произволни $x,y$:
\[
  \mathbb P(X<x,Y<y) = \sum_{x_i<x}\sum_{y_j<y}\mathbb P(X=x_i,Y=y_j) = \sum_{x_i<x}\sum_{y_j<y}\mathbb P(X=x_i)\mathbb P(Y=y_j)
\]
\[
  = \Big(\sum_{x_i<x}\mathbb P(X=x_i)\Big)\Big(\sum_{y_j<y}\mathbb P(Y=y_j)\Big) = F_X(x)F_Y(y).
\]

($\Rightarrow$) Нека $F_{X,Y}(x,y)=F_X(x)F_Y(y)$ за всички $(x,y)$.
Нека $x_i,y_j$ са съседни стойности в~носителите на $X,Y$, а~$x_i^-,y_j^-$ означават непосредствено предходната стойност (или
$-\infty$, ако $x_i$/$y_j$ е най-малката). По принципа на
включване-изключване:
\begin{align*}
  \mathbb P(X=x_i,Y=y_j) &= F_{X,Y}(x_i^+,y_j^+) - F_{X,Y}(x_i^{-,+},y_j^+) - F_{X,Y}(x_i^+,y_j^{-,+}) + F_{X,Y}(x_i^{-,+},y_j^{-,+})\\
  &= F_X(x_i)F_Y(y_j) - F_X(x_i^-)F_Y(y_j) - F_X(x_i)F_Y(y_j^-) + F_X(x_i^-)F_Y(y_j^-)\\
  &= \big[F_X(x_i)-F_X(x_i^-)\big]\big[F_Y(y_j)-F_Y(y_j^-)\big] = \mathbb P(X=x_i)\,\mathbb P(Y=y_j). \qedbg
\end{align*}
\end{proofbox}

\section{Ковариация}

\begin{defbox}[~(Ковариация)]
Нека $X,Y$ са случайни величини. Количеството
\[
  \operatorname{cov}(X,Y) := \mathbb E\big[(X-\mathbb EX)(Y-\mathbb EY)\big]
\]
се нарича \textbf{ковариация} на $X$ и~$Y$, ако съществува. Тя
измерва как двете случайни величини се изменят заедно. За дискретни
$X,Y$: $\operatorname{cov}(X,Y) = \sum_{i,j}(x_i-\mathbb EX)(y_j-\mathbb EY)\,p_{ij}$.
\end{defbox}

\begin{rembox}[Идея: линейна връзка]
Линейна зависимост между две случайни величини имаме, когато
$Y=aX+b$. Ковариацията е мярка за \emph{посоката} на линейната връзка.
Наистина, ако положим $\widetilde X=X-\mathbb EX$,
$\widetilde Y=Y-\mathbb EY$ (центрирани версии), то от $Y=aX+b$
следва
\[
  Y-\mathbb EY = a(X-\mathbb EX) \iff \widetilde Y = a\widetilde X,
\]
т.е.\ центрирането не променя линейната зависимост между величините
- то само премества центъра в~началото на координатната система.
\end{rembox}

\begin{thmbox}[]{Теорема}
Нека $X,Y$ са случайни величини. Тогава
\[
  \operatorname{cov}(X,Y) = \mathbb E\widetilde X\widetilde Y = \mathbb EXY - \mathbb EX\,\mathbb EY,
  \qquad \widetilde X=X-\mathbb EX,\ \widetilde Y = Y-\mathbb EY.
\]
\end{thmbox}

\begin{proofbox}
\begin{align*}
  \mathbb E(X-\mathbb EX)(Y-\mathbb EY) &= \mathbb E\big[XY - Y\,\mathbb EX - X\,\mathbb EY + \mathbb EX\,\mathbb EY\big]\\
  &= \mathbb EXY - \mathbb EX\,\mathbb EY - \mathbb EY\,\mathbb EX + \mathbb EX\,\mathbb EY = \mathbb EXY - \mathbb EX\,\mathbb EY. \qedbg
\end{align*}
\end{proofbox}

\begin{rembox}[Забележка: скалиране]
Преминаването $X,Y\to X-\mathbb EX,\,Y-\mathbb EY$ не променя
зависимостта между $X$ и~$Y$, а~само центрира величините. От друга
страна, скалирането $X,Y\to 10X,10Y$ скалира ковариацията
квадратично:
\[
  \operatorname{cov}(10X,10Y) = \mathbb E(10X-10\,\mathbb EX)(10Y-10\,\mathbb EY) = 100\,\mathbb E(X-\mathbb EX)(Y-\mathbb EY) = 100\operatorname{cov}(X,Y).
\]
\end{rembox}

\begin{thmbox}[~(на ковариацията)]{Свойства}
\begin{enumerate}[label=\arabic*)]
  \item \textbf{Симетрия:} $\operatorname{cov}(X,Y)=\operatorname{cov}(Y,X)$.
  \item $\operatorname{cov}(X,X)=DX$.
  \item \textbf{Линейност:} $\operatorname{cov}(aX+b,\,cY+d) = ac\cdot\operatorname{cov}(X,Y)$.
  \item \textbf{Адитивност:} $\operatorname{cov}(X+Z,\,Y) = \operatorname{cov}(X,Y)+\operatorname{cov}(Z,Y)$.
  \item Ако $X\perp\!\!\!\perp Y$, то $\operatorname{cov}(X,Y)=0$.
\end{enumerate}
\end{thmbox}

\begin{proofbox}
\textbf{1)} Директно от дефиницията: $\operatorname{cov}(X,Y)=\mathbb
E(X-\mathbb EX)(Y-\mathbb EY) = \mathbb E(Y-\mathbb EY)(X-\mathbb EX)
= \operatorname{cov}(Y,X)$.

\textbf{2)} $\operatorname{cov}(X,X) = \mathbb EX\cdot X - \mathbb
EX\cdot\mathbb EX = \mathbb EX^2-(\mathbb EX)^2 = DX$.

\textbf{3)} Ползвайки формулата $\operatorname{cov}(U,V)=\mathbb
EUV-\mathbb EU\,\mathbb EV$ и~линейността на очакването:
\begin{align*}
  \operatorname{cov}(aX+b,cY+d) &= \mathbb E\big[(aX+b)(cY+d)\big] - \mathbb E(aX+b)\,\mathbb E(cY+d)\\
  &= \mathbb E\big[ac\,XY+ad\,X+bc\,Y+bd\big] - (a\,\mathbb EX+b)(c\,\mathbb EY+d)\\
  &= ac\,\mathbb EXY+ad\,\mathbb EX+bc\,\mathbb EY+bd - \big(ac\,\mathbb EX\,\mathbb EY+ad\,\mathbb EX+bc\,\mathbb EY+bd\big)\\
  &= ac\big(\mathbb EXY-\mathbb EX\,\mathbb EY\big) = ac\cdot\operatorname{cov}(X,Y).
\end{align*}

\textbf{4)} Аналогично:
\begin{align*}
  \operatorname{cov}(X+Z,Y) &= \mathbb E\big[(X+Z)Y\big] - \mathbb E(X+Z)\,\mathbb EY = \mathbb EXY+\mathbb EZY - (\mathbb EX+\mathbb EZ)\mathbb EY\\
  &= \big(\mathbb EXY-\mathbb EX\,\mathbb EY\big) + \big(\mathbb EZY-\mathbb EZ\,\mathbb EY\big) = \operatorname{cov}(X,Y)+\operatorname{cov}(Z,Y).
\end{align*}

\textbf{5)} Ако $X\perp\!\!\!\perp Y$, по свойството за очакване на
произведение на независими случайни величини (Тема~5),
$\mathbb EXY=\mathbb EX\,\mathbb EY$, следователно
$\operatorname{cov}(X,Y)=\mathbb EXY-\mathbb EX\,\mathbb EY=0$.
\emph{(Обратното не е вярно в~общия случай - нулева ковариация не
означава непременно независимост.)} \qedbg
\end{proofbox}

\section{Корелация}

\begin{defbox}[~(Корелация)]
Нека $X,Y$ са случайни величини с~$0<DX<\infty$, $0<DY<\infty$
(т.е.\ нито едната не е почти сигурно константа). Тогава
\[
  \rho(X,Y) := \frac{\operatorname{cov}(X,Y)}{\sqrt{DX}\sqrt{DY}}
\]
се нарича \textbf{корелация} (коефициент на корелация) между $X$ и~$Y$.
\end{defbox}

\begin{rembox}[Инвариантност спрямо мащаб]
За разлика от ковариацията, корелацията е инвариантна спрямо
(положително) мащабиране:
\[
  \rho(10X,10Y) = \frac{100\operatorname{cov}(X,Y)}{\sqrt{D(10X)}\sqrt{D(10Y)}} = \frac{100\operatorname{cov}(X,Y)}{10\sqrt{DX}\cdot10\sqrt{DY}} = \frac{\operatorname{cov}(X,Y)}{\sqrt{DX}\sqrt{DY}} = \rho(X,Y).
\]
\end{rembox}

\begin{thmbox}[]{Теорема}
Нека $X,Y$ са случайни величини с~$0<DX<\infty$, $0<DY<\infty$, и~нека
$\bar X = \dfrac{X-\mathbb EX}{\sqrt{DX}}$, $\bar Y = \dfrac{Y-\mathbb
EY}{\sqrt{DY}}$ (\textbf{стандартизирани} случайни величини). Тогава
$\mathbb E\bar X = \mathbb E\bar Y = 0$, $D\bar X = D\bar Y = 1$, и~\[
  \rho(X,Y) = \mathbb E\bar X\bar Y.
\]
\end{thmbox}

\begin{proofbox}
$\mathbb E\bar X = \mathbb E\dfrac{X-\mathbb EX}{\sqrt{DX}} =
\dfrac1{\sqrt{DX}}\,\mathbb E(X-\mathbb EX) = 0$, понеже $\mathbb
E(X-\mathbb EX)=\mathbb EX-\mathbb EX=0$. Аналогично $\mathbb E\bar
Y=0$.

$D\bar X = D\Big(\dfrac{X-\mathbb EX}{\sqrt{DX}}\Big) =
\dfrac1{(\sqrt{DX})^2}\,D(X-\mathbb EX) = \dfrac1{DX}\cdot DX = 1$
(изваждането на константа $\mathbb EX$ не променя дисперсията).
Аналогично $D\bar Y=1$.

Накрая:
\[
  \rho(X,Y) = \frac{\operatorname{cov}(X,Y)}{\sqrt{DX}\sqrt{DY}} = \mathbb E\left(\frac{X-\mathbb EX}{\sqrt{DX}}\right)\left(\frac{Y-\mathbb EY}{\sqrt{DY}}\right) = \mathbb E\bar X\bar Y. \qedbg
\]
\end{proofbox}

\begin{thmbox}[]{Теорема}
Нека $X,Y$ са случайни величини с~$0<DX<\infty$, $0<DY<\infty$. Тогава:
\begin{enumerate}[label=\alph*)]
  \item $|\rho(X,Y)|\le 1$;
  \item $|\rho(X,Y)|=1 \iff Y=aX+b$ за някои $a,b\in\mathbb R$ (т.е.\ $Y$ се предопределя линейно от $X$).
\end{enumerate}
\end{thmbox}

\begin{proofbox}
\textbf{а)} Нека $\bar X,\bar Y$ са стандартизираните величини от
предната теорема, така че $\rho(X,Y)=\mathbb E\bar X\bar Y$. Понеже
$D\bar X = \mathbb E\bar X^2-(\mathbb E\bar X)^2=1$ и~$\mathbb
E\bar X=0$, следва $\mathbb E\bar X^2=1$; аналогично $\mathbb
E\bar Y^2=1$. Разглеждаме
\[
  0 \le \mathbb E(\bar X\pm\bar Y)^2 = \mathbb E\bar X^2 \pm 2\,\mathbb E\bar X\bar Y + \mathbb E\bar Y^2 = 1\pm2\rho(X,Y)+1 = 2\big(1\pm\rho(X,Y)\big),
\]
откъдето $1\pm\rho(X,Y)\ge0$ и~за двата знака, т.е.\ $-1\le\rho(X,Y)\le1$, или $|\rho(X,Y)|\le1$.

\textbf{б) ($\Rightarrow$)} Нека $\rho(X,Y)=1=\mathbb E\bar X\bar Y$.
От горното с~избор на знак „$-$“:
\[
  \mathbb E(\bar X-\bar Y)^2 = 1-2\rho(X,Y)+1 = 1-2+1 = 0.
\]
Понеже $(\bar X-\bar Y)^2\ge0$ и~очакването му е нула, следва
$(\bar X-\bar Y)^2=0$ (с~вероятност~1), т.е.\ $\bar X=\bar Y$. Записвайки
explicitно $\bar X = \dfrac{X-\mathbb EX}{\sqrt{DX}}$, $\bar Y =
\dfrac{Y-\mathbb EY}{\sqrt{DY}}$, от $\bar X=\bar Y$ получаваме
\[
  Y = \sqrt{\frac{DY}{DX}}\,X + \Big(\mathbb EY - \sqrt{\frac{DY}{DX}}\,\mathbb EX\Big) =: aX+b,
\]
където $a=\sqrt{DY/DX}\ge0$ (в~случая $\rho(X,Y)=-1$ разсъждението е
аналогично с~избор на знак „$+$“, и~се получава $a\le0$).

\textbf{б) ($\Leftarrow$)} Нека $Y=aX+b$. Тогава $\mathbb EY=a\,
\mathbb EX+b$, $DY=D(aX+b)=a^2DX$, и~\begin{align*}
  \rho(X,Y) &= \frac{\mathbb E(X-\mathbb EX)(Y-\mathbb EY)}{\sqrt{DX}\sqrt{DY}} = \frac{\mathbb E(X-\mathbb EX)(aX+b-a\,\mathbb EX-b)}{\sqrt{DX}\sqrt{a^2DX}}\\
  &= \frac{a\,\mathbb E(X-\mathbb EX)^2}{|a|\,DX} = \frac{a\,DX}{|a|\,DX} = \frac{a}{|a|} = \pm1. \qedbg
\end{align*}
\end{proofbox}
